在 \([0,1]\) 上取 \(f_n(x)=n\,\mathbf1_{(0,1/n]}(x)\)。固定任何一个 \(x>0\)，只要 \(n\) 够大就有 \(f_n(x)=0\)，所以逐点极限处处为零，\(\int\lim_n f_n=0\)。而每一个 \(f_n\) 的积分都等于 1。一个恒等于零的极限函数和一列积分恒为 1 的函数，同时成立。

这个反例把本章要处理的困难摆在了明处：逐点收敛描述的是每个位置上的读数，而积分统计的是总量，二者之间没有自动的桥梁。质量可以在极限过程中挤成尖峰或者漂向远方，逐点的信息对此一无所知。凡是想把极限号、求和号或求导号推进积分号的地方，都要先补上一个额外条件，说明质量跑不掉。

\begin{center}
\begin{tikzpicture}[>=stealth,scale=0.85]
\draw[->] (0,0) -- (4.6,0) node[right] {\footnotesize \(x\)};
\draw[->] (0,0) -- (0,3.2);
\draw[thick,dotted] (0,0.7) -- (4,0.7) -- (4,0);
\draw[thick,dashed] (0,1.4) -- (2,1.4) -- (2,0);
\draw[thick] (0,2.8) -- (1,2.8) -- (1,0);
\node at (2.2,-0.75) {\footnotesize 面积恒为 \(1\)};
\begin{scope}[shift={(6.6,0)}]
\draw[->] (0,0) -- (4.6,0) node[right] {\footnotesize \(x\)};
\draw[->] (0,0) -- (0,3.2);
\draw[very thick] (0,0) -- (4,0);
\node[right] at (1.2,0.5) {\footnotesize \(\lim_n f_n\equiv 0\)};
\node at (2.2,-0.75) {\footnotesize 面积等于 \(0\)};
\end{scope}
\end{tikzpicture}
\end{center}

\emph{左右两张图之间少了一个条件，本章要找的就是这个条件。}

本章先让交换失败，再逐步补上足以拦住失败的条件：

\begin{enumerate}
\tightlist
\item
  点态极限为何不能直接穿过积分号：用移动尖峰说明逐点收敛不约束总质量，再给出一致收敛、单调收敛与 Fatou 引理三种补救。
\item
  控制收敛用一个可积包络守住质量：一个与 \(n\) 无关的可积上界同时封死``长成尖峰''与``漂向远方''两条通道，并把结论平移到期望与无穷级数上。
\item
  求导穿过积分号需要控制差商：被控制的对象换成差商，条件的形式不变；似然求导、重参数化梯度与 score function 估计都建立在这一步上。
\end{enumerate}

单调收敛与控制收敛都是 Lebesgue 积分的定理，需要可测性与相应的可积条件。只学过 Riemann 积分的读者可以先在有限区间上的连续函数这一特例里掌握它们的用法，但不要把理论框架略过去——尖峰反例恰恰是在 Riemann 视角下最容易被当成``显然可以交换''的那一类。

\subsection{一个具体算例}

仍用开头的 \(f_n=n\mathbf1_{(0,1/n]}\)。Fatou 引理给出 \(\int\liminf_n f_n\le\liminf_n\int f_n\)，此处是 \(0\le1\)，不等式成立且不等号严格。丢掉的那一份就是逃走的质量。追问一句为什么控制收敛定理救不了：逐点上确界 \(\sup_n f_n(x)\) 在零点附近与 \(1/x\) 同阶，积分发散，整族头顶没有可积的天花板。

\subsection{怎么读这一章}

三篇围绕同一个句式展开：什么时候可以交换两个操作。读的时候每次都把两件事说清楚——哪一列函数在收敛，以及谁在提供与序号无关的统一控制。第三篇里被交换的对象从函数变成了差商，如果前两篇的这两个问题回答得清楚，第三篇几乎是照搬。

\subsection{本章篇目}

以下篇目按概念依赖排列；若从具体问题进入，也可以先打开最接近当前困惑的一篇，再沿篇内链接回补。

\begin{itemize}
\tightlist
\item \hyperref[sec:12-Real-Analysis-Support/04-Limits-and-Integration/01]{``点态极限为何不能直接穿过积分号''}；
\item \hyperref[sec:12-Real-Analysis-Support/04-Limits-and-Integration/02]{``控制收敛用一个可积包络守住质量''}；
\item \hyperref[sec:12-Real-Analysis-Support/04-Limits-and-Integration/03]{``求导穿过积分号需要控制差商''}。
\end{itemize}

读完三篇，凡是写下``交换极限与积分''的地方，你都能立刻说出所依据的那个统一控制量是谁。
